\begin{abstract}
Interactive imitation learning lets a learner practise a task and an expert take
over to demonstrate the right actions. Existing methods decide one thing: when to
take over. They run an episode from its start to its end, and they call the expert
at the first step where the policy becomes uncertain. Two decisions are left to
chance. The first is which failed episode to correct, where a failed episode is one
that does not achieve the task. The second is where the demonstration begins. Each
demonstration costs human time, and only so many can be collected, so both
decisions govern what that effort buys. We present \textbf{DISEIL}
(\textbf{D}emonstration d\textbf{I}stillation for \textbf{S}ample-\textbf{E}fficient
\textbf{I}mitation \textbf{L}earning), which decides both. DISEIL sorts a round's
failed episodes into modes of failure. It then prescribes where a demonstration not
yet collected should begin, and checks that this starting configuration is
reachable before the expert is called. We evaluate across five tasks, with image
and state observations, ten settings in all, at a budget of 20 demonstrations,
against the best baselines. DISEIL reaches the highest mean success rate in every
setting, on average about 3 percentage points above the strongest baseline.
\end{abstract}
